\section{Conclusion}
\label{sec:conclusion}

We have presented \bench{}, a controlled micro-benchmark for evaluating VLM robustness to UI layout and style perturbations under grounded retrieval conditions. Through systematic evaluation of three frontier models across 300 images, 390 QA pairs, and five drift types at calibrated severity levels, we have identified several findings with practical consequences. Drift robustness is decoupled from base accuracy, with Gemini~2.0~Flash achieving the best DRS (0.950) despite lower base EM than GPT-4o. Evidence grounding remains critically weak across all models, with IoU below 0.07 even in the best case. Task-specific vulnerability profiles differ substantially across models, with Gemini failing on trend detection and GPT-4o-mini failing on table lookup, suggesting that model selection for UI agents should be guided by the expected task distribution.

\paragraph{Limitations.}
We acknowledge several limitations of this work. First, our benchmark uses synthetically generated UI pages rather than screenshots captured from live websites. While synthetic generation enables precise control over ground-truth annotations, it may not capture the full visual complexity of real-world web applications, including advertisements, dynamic content loading, and cross-browser rendering differences. Second, our evaluation covers three models from two providers; a more comprehensive evaluation would include open-source VLMs (LLaVA~\cite{liu2023llava}, Qwen-VL~\cite{bai2023qwenvl}), on-device models, and additional commercial systems. Third, all content is in English, and UI drift effects may interact differently with non-Latin scripts, right-to-left layouts, or multilingual interfaces. Fourth, the benchmark contains 50 base pages and 390 QA pairs; while sufficient for identifying model-level trends, this scale limits the statistical power of fine-grained per-type analyses.

\paragraph{Future work.}
Several directions extend naturally from this benchmark. Real-world web capture, using automated browsing pipelines to screenshot live websites under controlled viewport and theme configurations, would bridge the gap between synthetic and naturalistic evaluation. Expanding the model set to include open-source VLMs and specialized UI models (ScreenAI, Ferret-UI, CogAgent) would provide a more complete robustness landscape. Cross-lingual evaluation with non-English UIs and right-to-left scripts would assess whether drift robustness interacts with language-specific visual patterns. Adversarial drift, where perturbations are specifically optimized to degrade model performance, would provide upper-bound degradation estimates. Finally, investigating training-time interventions (drift augmentation, consistency regularization) that improve robustness without sacrificing base accuracy represents a high-impact direction for making VLM agents reliable in production UI settings.

\bench{} provides a systematic foundation for measuring a dimension of VLM reliability that existing benchmarks overlook. As VLM agents are deployed in increasingly consequential UI automation tasks, ensuring robust performance under the visual variability inherent to real-world interfaces is not merely a desirable property but a prerequisite for trustworthy deployment.
